% ============================================================
%  CP-Lib 打印版 —— 导言区
%  引擎：XeLaTeX（中文必需）。用 latexmk -xelatex 编译。
% ============================================================

% ---------- 页面：A4 纵向，双栏，窄边距 ----------
% top 必须容纳 headheight + headsep，否则页眉会被裁掉
\usepackage[a4paper,top=13mm,bottom=12mm,left=9mm,right=9mm,headheight=13pt,headsep=3.5mm]{geometry}
\usepackage{multicol}

% ---------- 字体 ----------
% ctexart 已自动加载 xeCJK 并选用 Windows 字体（SimSun/SimHei/KaiTi）。
% 代码块里的中文注释走 \CJKttfamily，这里显式指定为宋体，小字号下最清晰。
\setCJKmonofont{SimSun}
% 拉丁等宽字体用于代码；Consolas 是 Windows 自带、比赛环境常见
\setmonofont{Consolas}[Scale=0.92]

% ---------- 数学 ----------
\usepackage{amsmath}
\usepackage{amssymb}
% CP-Lib.md 第 3565 行用了 $180\degree$，而 \degree 并非 base LaTeX/amsmath 提供，
% 不定义会直接编译失败。这里补上。
\newcommand{\degree}{\ensuremath{^{\circ}}}

% ---------- 标题层级压缩（双栏下必须，否则标题吃掉大量版面）----------
\usepackage{titlesec}
\usepackage{titletoc}

% h1 -> \section
\titleformat{\section}
  {\normalfont\large\bfseries\color{seccolor}}
  {\thesection}{0.6em}{}
\titlespacing*{\section}{0pt}{7pt}{3pt}

% h2 -> \subsection
\titleformat{\subsection}
  {\normalfont\normalsize\bfseries\color{subseccolor}}
  {\thesubsection}{0.5em}{}
\titlespacing*{\subsection}{0pt}{5pt}{2pt}

% h3 -> \subsubsection
\titleformat{\subsubsection}
  {\normalfont\small\bfseries}
  {\thesubsubsection}{0.5em}{}
\titlespacing*{\subsubsection}{0pt}{4pt}{1.5pt}

% h4 -> \paragraph（不编号，紧凑）
\titleformat{\paragraph}[runin]
  {\normalfont\small\bfseries}
  {}{0em}{}
\titlespacing*{\paragraph}{0pt}{3pt}{0.5em}

% 目录：收 h1/h2/h3
\setcounter{tocdepth}{3}
\setcounter{secnumdepth}{3}
\contentsmargin{1.6em}
\titlecontents{section}[0em]{\vspace{1pt}\bfseries}{\thecontentslabel\quad}{}{\titlerule*[0.4pc]{.}\contentspage}
\titlecontents{subsection}[1.2em]{\vspace{0.5pt}}{\thecontentslabel\quad}{}{\titlerule*[0.4pc]{.}\contentspage}
\titlecontents{subsubsection}[2.9em]{\small}{\thecontentslabel\quad}{}{\titlerule*[0.4pc]{.}\contentspage}

% ---------- 颜色 ----------
\usepackage{xcolor}
\definecolor{seccolor}{RGB}{0,0,0}
\definecolor{subseccolor}{RGB}{20,20,20}
\definecolor{codekw}{RGB}{0,0,190}
\definecolor{codecmt}{RGB}{0,120,0}
\definecolor{codestr}{RGB}{170,20,20}
\definecolor{codenum}{RGB}{140,0,140}
\definecolor{coderule}{RGB}{190,190,190}
\definecolor{linkcolor}{RGB}{0,0,200}

% ---------- 代码高亮 ----------
\usepackage{listings}
\lstset{
  language=C++,
  basicstyle=\ttfamily\scriptsize,
  keywordstyle=\color{codekw}\bfseries,
  commentstyle=\color{codecmt}\itshape,
  stringstyle=\color{codestr},
  numberstyle=\tiny\color{gray},
  % 折行：代码行宽 p99 约 85 列，单栏约容纳 69 列，超长行必须折
  breaklines=true,
  breakatwhitespace=false,
  breakindent=1em,
  % 注意：prebreak/postbreak 位于 \discretionary 内，其中不能出现 whatsit
  % （\color 会产生 whatsit），否则报 "Improper discretionary list"。故用 \raisebox。
  prebreak=\raisebox{0ex}[0ex][0ex]{\ensuremath{\searrow}},
  postbreak=\raisebox{0ex}[0ex][0ex]{\ensuremath{\hookrightarrow\space}},
  columns=fullflexible,
  keepspaces=true,
  showstringspaces=false,
  tabsize=4,
  extendedchars=true,
  frame=l,
  framerule=0.4pt,
  rulecolor=\color{coderule},
  xleftmargin=1.4mm,
  xrightmargin=0.4mm,
  framexleftmargin=0.4mm,
  aboveskip=4pt,
  belowskip=3pt,
  % 超长代码块允许跨栏/跨页断开
  breakautoindent=true,
}

% ---------- 页眉页脚（页码）----------
\usepackage{fancyhdr}
\usepackage{lastpage}
\pagestyle{fancy}
\fancyhf{}
\renewcommand{\headrulewidth}{0.3pt}
\renewcommand{\footrulewidth}{0pt}
\renewcommand{\sectionmark}[1]{\markboth{#1}{}}
\fancyhead[L]{\small\nouppercase{\leftmark}}
\fancyhead[R]{\small\itshape CP-Lib}
\fancyfoot[C]{\small\thepage\ /\ \pageref{LastPage}}

% ---------- 图片、链接 ----------
\usepackage{graphicx}
\usepackage[hidelinks]{hyperref}
\hypersetup{
  colorlinks=true,
  linkcolor=linkcolor,
  urlcolor=linkcolor,
  bookmarksnumbered=true,
  pdftitle={CP-Lib 算法竞赛模板},
  pdfsubject={算法竞赛模板库},
}

% ---------- 列表 ----------
\usepackage{enumitem}
\setlist[itemize]{topsep=2pt,itemsep=1pt,parsep=0pt,leftmargin=1.2em}
\setlist[enumerate]{topsep=2pt,itemsep=1pt,parsep=0pt,leftmargin=1.6em}

% ---------- 正文排版 ----------
\usepackage{indentfirst}
\setlength{\parindent}{0pt}
\setlength{\parskip}{2pt}
\linespread{1.04}
